\section{Chapters 20 and 21}

\label{sec:PrincipleAction2021}

\paragraph{Chapter 20}

\begin{verse}
The (discursive, reasoning) intellect is the source of sorrow\\ What consistency do dialectical distinctions have?\\ What is the difference between ``good'' and ``evil''?\\ ``To act like everyone else'' --- a norm based on fear.\\ No! (out of) this squalor.

People are carried away by a simple joy\\ A festival suffices for them\\ A panorama of spring which is seen from a terrace\\ I, instead, am anchored in the bottom of the current of sentiment\\ I remain serious and calm beyond joy like an infant who has not yet smiled\\ I live like this and go as if I belong to no place.\\ Everyone desires excess\\ While I am like one who possesses nothing\\ (I seem) ignorant, simple, without a practical spirit\\ Everyone lives in the light\\ While I am darkness\\ Everyone needs company\\ While I love only the solitary height\\ Indeterminable like the waters of the ocean\\ Like the wave I turn without pause\\ Everyone aspires to something\\ While I seem simple and incompetent\\ I am different from everyone\\ But, united to the original productive essence, only I am an I.
\end{verse}

\paragraph{Chapter 21}

\begin{verse}
The great Virtue which is manifested\\ Is only the exteriorization of the Principle\\ But the substance of the Principle\\ Is undifferentiated and unintelligible\\ Undifferentiated and unintelligible\\ It contains the seeds the archetypes, the seminal possibilities of beings, constituting their indestructible element, their ``heaven'' \\ Mysterious and incomprehensible\\ It contains existences = beings who have passed to the formal state \\ Deep and hidden\\ It contains essences transformed beings --- ``spirits'' \\ As such it is the great Reality\\ The secure place (of all beings).\\ From the origin up to the present\\ It does not change its Name\\ The animating principle (of each thing) proceeds from it\\ What is the foundation of such knowledge?\\ This.
\end{verse}

\paragraph{Commentary}


Beings do not cease to live in the Principle and to have in it their ``secure place'' in the triple state:

\begin{itemize}


\item In the preformal state ``being part of Heaven and Earth'' as archetypes or seeds (\emph{logoi spermatikoi})

\item In formal existence or existence in its proper and limitative sense

\item As pure forms.
\end{itemize}


Coexistence of the emission, of the eternal movement where those states and the mysterious, unintelligible, unmoved Unity follow each other.

To the unknowability and the unintelligibility which would be in the same way the predicate of the Principle is contrasted, in the last lines, with the form proper to a possible knowledge of its mystery: ``this'', the word refers to direct, metaphysical, and super-rational intuition (it appears again in Chapter 54). The expression \emph{tathata} may be compared, which, when used with an analogous meaning in \textbf{Madhyamika Buddhism}, designated the content of the supreme, inexpressable experience.

\paragraph{Chinese text and literal translation}


Chapter 20 (\begin{CJK*}{UTF8}{bsmi}第二十章\end{CJK*})
\columnratio{0.37}
\begin{paracol}{2}
\begin{CJK*}{UTF8}{bsmi}
\begin{verse}
絕學無憂。\\ 唯之與阿，相去幾何？\\ 善之與惡，相去何若？\\ 人之所畏，不可不畏？\\ 荒兮其未央哉！\\ 眾人熙熙，\\ 如享太牢，\\ 如登春臺。\\ 我獨泊兮其未兆，\\ 若嬰兒之未孩，\\ 乘乘兮若無所歸。\\ 眾人皆有餘，\\ 而我獨若遺。\\ 我愚人之心也哉！沌沌兮！\\ 俗人昭昭，我獨昏昏；\\ 俗人察察，我獨悶悶。\\ 忽兮其若晦，\\ 寂兮似無所止。\\ 眾人皆有以，我獨頑且鄙。\\ 我獨異於人，\\ 而貴求食於母。
\end{verse}
\end{CJK*}
\switchcolumn
\vspace*{-2ex}
\begin{verse}
Discard conventional doctrines and be relieved from anxieties.\\ Flattery or reprimand, what difference does it make?\\ Good or evil, what difference does it take?\\ Just because the people are at awe, you cannot be indifferent?\\ Ridiculous! Ungrounded!\\ When everyone is celebrating and joyous,\\ As if relishing a spiritual triumph,\\ As if enjoying a victorious feast.\\ I alone am feelingless! Contemplating the myriad future,\\ Dazed like a newborn child,\\ Seizing the moment! Pondering the uncertainties.\\ When everyone feels lavished,\\ I alone am hollowed.\\ I am a fool! Confounded!\\ When everyone seemed enlightened, I alone am in suspense;\\ When everyone is watchful, I alone am in confuse.\\ Enigmatic! Like the obscure twilight,\\ Desolate! Like the infinite universe.\\ When everyone is focused, I alone am stubborn and despicable.\\ I alone am different from the common,\\ I find refuge in the womb of this profound conception.
\end{verse}
\end{paracol}

Chapter 21 (\begin{CJK*}{UTF8}{bsmi}第二十一章\end{CJK*})
\begin{paracol}{2}
\begin{CJK*}{UTF8}{bsmi}
\begin{verse}
孔德之容，\\ 唯道是從。\\ 道之為物，\\ 惟恍惟惚。\\ 惚兮恍其中有象；\\ 恍兮惚其中有物。\\ 窈兮冥兮，\\ 其中有精；\\ 其精甚真，其中有信。\\ 自古及今，\\ 其名不去，\\ 以閱眾甫。\\ 吾何以知眾甫之然哉？\\ 以此。
\end{verse}
\end{CJK*}
\switchcolumn
\vspace*{-2ex}
\begin{verse}
Where the greatest Virtue resides,\\ Only the Dao may reveal.\\ Things that bestow the Dao,\\ Radiates such liberty, such easiness.\\ Eased! Liberated from its form, yet fashioned;\\ Liberated! At ease with its position, yet poised.\\ Mesmerizing! Mysterious!\\ Brilliance from within;\\ Its brilliance so enlightening, it exemplifies the truth.\\ Traversing through all times,\\ Its name is not diminished,\\ Amassing all marvels of human intelligence.\\ How do I know the essence of all marvels?\\ By observing things that bestow the Dao.
\end{verse}
\end{paracol}


\flrightit{Posted on 2021-03-19 by Lao Tzu}

